\section{Literature Review}
\label{sec:litreview}

\subsection{Identification of Common Genetic Risk Variants for Autism Spectrum Disorder}

Grove et al.\ (2019)~\cite{grove2019}, published in \textit{Nature
Genetics}, represents the largest genetic study of autism using common
variants to date. By combining data from over 46\,000 individuals, the
study confirmed that common genetic variants collectively play a
meaningful role in autism risk and identified specific regions of the
genome where this risk is concentrated. However, when these findings
were used to build a prediction tool, it explained only 2.5\,\% of why
some people develop autism and others do not, despite the enormous
sample size. The authors concluded that better prediction would require
fundamentally different approaches rather than simply collecting more
data, leaving a clear methodological gap that subsequent work has yet
to fill.

\subsection{Large-Scale Exome Sequencing Study Implicates Both Developmental and Functional Changes in the Neurobiology of Autism}

Satterstrom et al.\ (2020)~\cite{satterstrom2020}, published in
\textit{Cell}, offers an important clue as to what those different
approaches might target. By examining rare mutations in 35\,584
individuals, the study identified 102 genes with strong links to
autism. These genes were not randomly distributed across biology but
clustered around two specific functions: controlling brain development
before birth and managing communication between neurons. More
importantly, these same genes overlapped with the regions flagged by
common variant studies like Grove et al., suggesting that different
types of genetic risk ultimately disrupt the same underlying
biological systems. Autism risk appears to funnel through a connected
set of biological pathways rather than arising from isolated genetic
effects acting independently, which points toward gene interaction
modelling as a more biologically grounded basis for prediction than
current additive methods provide.

\subsection{Understanding the Genetic Contribution to ASD}

Havdahl et al.\ (2021)~\cite{havdahl2021} review synthesises the
current understanding of ASD's genetic architecture, emphasising that
risk arises from two distinct sources. Rare \textit{de novo} mutations
in constrained genes like \textit{CHD8} and \textit{SCN2A} confer large
individual effects but collectively explain fewer than 5\,\% of
cases. Common inherited variants each have tiny effects, and despite
large GWAS efforts, only a handful of loci have been robustly
identified. SNP-based heritability estimates range widely from 12\,\%
to 65\,\%, far below the 80\,\% from twin studies. The authors argue
that closing this gap requires integrating common and rare variant
data alongside multi-omics approaches, as standard additive models are
insufficient to capture the full genetic liability. This directly
motivates the need for models that go beyond simple summation of
variant effects and instead account for how genes relate to and
interact with one another.

\subsection{Single Nucleotide Polymorphisms Predict Symptom Severity of Autism Spectrum Disorder}

Jiao et al.\ (2012)~\cite{jiao2012} investigated whether the genetic
profile of a child with autism, specifically their pattern of SNPs
across a set of known autism-related genes, can predict how severe
their symptoms will be. Working with 118 children already diagnosed
with ASD, the researchers used a standard clinical severity scale
called the Childhood Autism Rating Scale (CARS) to divide participants
into two groups: those with mild to moderate autism and those with
severe autism. For each child, 29 SNPs spanning 9 genes previously
linked to ASD were collected and fed into three machine learning models
-- decision stumps, alternating decision trees, and FlexTrees -- to
see whether genetic data alone could correctly classify symptom
severity. The best-performing models achieved an overall accuracy of
67\,\%, with a sensitivity of 0.88, meaning they were reasonably good
at identifying the more severely affected children. Across all three
models, one SNP consistently emerged as the most important predictor:
rs878960, located in a gene called \textit{GABRB3}, which plays a role
in how the brain regulates key chemical signals between neurons.
Children carrying certain versions of this SNP tended to score higher
on the autism severity scale, suggesting a genuine biological link
between this genetic variant and the intensity of autistic
symptoms. This paper is significant for our work for two reasons.
First, it provides direct evidence that SNP-level genetic data carries
real predictive signal about autism --- not just who develops the
condition, but how it manifests --- which directly supports the
premise of our project. Second, it honestly reveals the limitations of
a small, focused approach: only 118 children, only 29 SNPs from
pre-selected genes, and modest accuracy. Our research aims to extend
this logic to a genome-wide scale, using thousands of SNPs and more
powerful machine learning methods to substantially improve on what
this foundational study was able to achieve.

\subsection{Interactions of Genetic Risks for Autism and the Broad Autism Phenotypes}

Dong et al.\ (2023)~\cite{dong2023} takes a bioinformatics pipeline
approach to combining SNP-level GWAS data with machine
learning-informed polygenic risk scoring, making it highly relevant to
the kind of work our team is pursuing. The researchers worked with
genome-wide SNP data from nearly 10\,000 individuals including over
2\,400 autism cases and their unaffected siblings, computing polygenic
risk scores using a Bayesian method called PRS-CS that estimates the
effect of each SNP while accounting for correlations between nearby
variants. GWAS summary statistics from a study involving over 18\,000
cases and 27\,000 controls were combined with the European 1000
Genomes reference panel to estimate posterior effect sizes, generating
polygenic scores across 3\,659\,780 SNPs for nearly 10\,000
individuals. What makes this paper particularly valuable is that it
goes beyond simply computing a risk score and instead investigates how
that SNP-derived genetic risk interacts with other factors,
specifically rare \textit{de novo} mutations, to shape autism outcomes
and related behavioural traits across the spectrum. The study found
that common SNP-based polygenic risk and rare mutations contribute to
autism risk through overlapping but distinct pathways, and that their
combined effect is greater than either alone. For our project, this
paper serves as a strong example of how a well-designed bioinformatics
pipeline that properly handles the full SNP dataset rather than only
keeping a few significant variants can reveal biological relationships
that simpler approaches would completely miss.

\subsection{Problem Statement}

Genome-Wide Association Studies (GWAS) have served as the primary
scientific tool for uncovering the genetic foundations of Autism
Spectrum Disorder. By scanning the genomes of thousands of
individuals, some diagnosed with ASD and others without, researchers
can identify specific positions in the DNA, known as single nucleotide
polymorphisms or SNPs, where particular genetic variants appear more
frequently among those with autism. This method has illuminated a
number of genomic regions influencing the neurobiological processes
underlying ASD, offering meaningful, if incomplete, biological insight.

To translate these findings into something clinically actionable,
researchers developed the Polygenic Risk Score, or PRS, a tool that
combines the weighted contributions of many SNPs into a single number
representing an individual's genetic predisposition to ASD. In
principle, this is precisely the right direction, because autism risk
is not driven by any single variant but by the collective, cumulative
influence of potentially hundreds or thousands of SNPs acting together
across the genome. The critical problem, however, is that the majority
of existing bioinformatics research on ASD risk has not meaningfully
embraced this multi-SNP reality. Most studies either focus on
identifying individual SNPs of statistical significance, examine small
pre-selected sets of variants from a handful of candidate genes, or
apply machine learning to narrow genetic features rather than to the
full breadth of genome-wide SNP data. The rich, distributed signal
embedded across the entire SNP landscape therefore remains largely
unexploited.

The standard construction method known as pruning and thresholding
compounds this problem further by actively discarding the majority of
SNPs that do not individually clear a stringent significance
threshold. For a highly polygenic condition like ASD, where risk
emerges from the combined effect of potentially thousands of
small-effect variants, this means eliminating most of the true genetic
signal before any modelling even begins. Studies such as Jiao et al.\
(2012)~\cite{jiao2012}, while demonstrating that SNP patterns carry
genuine predictive information about autism, worked with only 29 SNPs
from 9 pre-selected genes, leaving the question of what a genuinely
genome-wide multi-SNP approach could achieve entirely unanswered. More
recent polygenic scoring work has moved closer to this goal but
continues to rely on additive linear frameworks that treat each SNP as
independent, ignoring the biological relationships between variants
that likely account for a substantial share of unexplained heritability.

There is therefore a clear and unresolved gap in the existing
literature: no bioinformatics approach for ASD has yet constructed a
risk prediction model that simultaneously incorporates the full
multi-SNP signal from genome-wide summary statistics, applies machine
learning to intelligently select and weight across that complete SNP
space, and produces a validated risk score with demonstrably superior
discriminatory accuracy. This is the precise gap the proposed research
sets out to address.
